\documentclass[11pt,a4paper]{article}
\usepackage[utf8]{inputenc}
\usepackage[T1]{fontenc}
\usepackage[spanish]{babel}
\usepackage{amsmath,amssymb,mathtools}
\usepackage{geometry}
\usepackage{enumitem}
\usepackage{array}
\usepackage{lmodern}
\usepackage{microtype}
\geometry{margin=2.4cm}
\setlength{\parindent}{0pt}
\setlength{\parskip}{0.45em}

\newcommand{\R}{\mathbb{R}}
\newcommand{\C}{\mathbb{C}}
\newcommand{\K}{\mathbb{K}}
\newcommand{\Ker}{\operatorname{Ker}}
\newcommand{\Ima}{\operatorname{Im}}
\newcommand{\rg}{\operatorname{rg}}
\newcommand{\tr}{\operatorname{tr}}
\newcommand{\Id}{\operatorname{Id}}
\newcommand{\diag}{\operatorname{diag}}
\newcommand{\Span}{\operatorname{L}}

\begin{document}

\begin{center}
{\LARGE\bfseries Álgebra Lineal}\\[0.4em]
{\Large Tema 3: Endomorfismos. Diagonalización --- Soluciones}
\end{center}

\section{Ejercicio 1}
\subsection{Solución}

Como $8$ es autovalor:

\[
\det(A-8I)=0
\]

por lo que:

\[
\rg(A-8I)<4.
\]

Además, $A-8I$ no puede tener rango $4$ y:

\[
\dim S(8)=4-\rg(A-8I).
\]

Por tanto:

\[
\boxed{
\rg(A-8I)\in\{0,1,2,3\}
}
\]

y:

\[
\boxed{
\dim S(8)\in\{1,2,3,4\}.
}
\]

\section{Ejercicio 2}
\subsection{Solución}

Tomamos:

\[
P=
\begin{pmatrix}
4&1\\
1&1
\end{pmatrix},
\qquad
D=
\begin{pmatrix}
3&0\\
0&0
\end{pmatrix}.
\]

Entonces:

\[
M=PDP^{-1}.
\]

Como:

\[
P^{-1}
=
\frac13
\begin{pmatrix}
1&-1\\
-1&4
\end{pmatrix},
\]

se obtiene:

\[
\boxed{
M=
\begin{pmatrix}
4&-4\\
1&-1
\end{pmatrix}.
}
\]

\section{Ejercicio 3}
\subsection{Solución}

Los autovalores de $A$ son:

\[
1,-1,2.
\]

Por las propiedades de los autovalores:

\[
\boxed{
4A:\quad 4,-4,8
}
\]

\[
\boxed{
A^t:\quad 1,-1,2
}
\]

\[
\boxed{
A^3:\quad 1,-1,8
}
\]

\[
\boxed{
A^{-1}:\quad 1,-1,\frac12.
}
\]

\section{Ejercicio 4}
\subsection{Solución}

Una matriz con todos sus autovalores simples es diagonalizable.

Por tanto:

\[
\boxed{
\text{La matriz }2\times2\text{ con autovalores }3,5\text{ es diagonalizable.}
}
\]

\[
\boxed{
\text{La matriz }3\times3\text{ con autovalores }6,7,8\text{ es diagonalizable.}
}
\]

Si existe un autovalor doble, hace falta conocer su multiplicidad geométrica.

Por tanto:

\[
\boxed{
\text{Con un autovalor }-3\text{ doble no hay información suficiente.}
}
\]

y:

\[
\boxed{
\text{Con }9\text{ doble y }1\text{ simple es diagonalizable}
\iff
\dim S(9)=2.
}
\]

\section{Ejercicio 5}
\subsection{Solución}

Sea:

\[
B=\{v_1,\ldots,v_n\}
\]

una base formada por autovectores.

Entonces existen:

\[
\lambda_1,\ldots,\lambda_n
\]

tales que:

\[
f(v_i)=\lambda_i v_i.
\]

Las coordenadas de $f(v_i)$ en $B$ son:

\[
[f(v_i)]_B
=
(0,\ldots,0,\lambda_i,0,\ldots,0)^t.
\]

Por tanto:

\[
\boxed{
M_B(f)
=
\begin{pmatrix}
\lambda_1&0&\cdots&0\\
0&\lambda_2&\cdots&0\\
\vdots&\vdots&\ddots&\vdots\\
0&0&\cdots&\lambda_n
\end{pmatrix}.
}
\]

\section{Ejercicio 6}
\subsection{Solución}

Para:

\[
A=
\begin{pmatrix}
a&b\\
c&d
\end{pmatrix},
\]

\[
P(\lambda)
=
\det(A-\lambda I)
=
(a-\lambda)(d-\lambda)-bc.
\]

Así:

\[
\boxed{
P(\lambda)
=
\lambda^2-(a+d)\lambda+ad-bc
=
\lambda^2-(\tr A)\lambda+\det A.
}
\]

Para una matriz $3\times3$:

\[
A=
\begin{pmatrix}
a&b&c\\
d&e&f\\
g&h&i
\end{pmatrix},
\]

desarrollando:

\[
\boxed{
P(\lambda)
=
-\lambda^3
+
(\tr A)\lambda^2
-
(A_{11}+A_{22}+A_{33})\lambda
+
\det A.
}
\]

\section{Ejercicio 7}
\subsection{Solución}

El único autovalor es:

\[
\lambda=a
\]

con multiplicidad algebraica $n$.

$A$ es diagonalizable si y sólo si:

\[
\dim\Ker(A-aI)=n.
\]

Esto equivale a:

\[
\rg(A-aI)=0.
\]

Por tanto:

\[
A-aI=0.
\]

Así:

\[
\boxed{
A\text{ es diagonalizable}
\iff
A=aI_n.
}
\]

\section{Ejercicio 8}
\subsection{Solución}

\[
\det(A-\lambda I)
=
\det((A-\lambda I)^t)
=
\det(A^t-\lambda I).
\]

Por tanto:

\[
\boxed{
A\text{ y }A^t\text{ tienen los mismos autovalores.}
}
\]

No tienen necesariamente los mismos autovectores. Por ejemplo:

\[
A=
\begin{pmatrix}
0&1\\
0&0
\end{pmatrix},
\qquad
A^t=
\begin{pmatrix}
0&0\\
1&0
\end{pmatrix}.
\]

Para $\lambda=0$:

\[
S_A(0)=\Span\{(1,0)\}
\]

mientras que:

\[
S_{A^t}(0)=\Span\{(0,1)\}.
\]

\[
\boxed{
\text{En general no tienen los mismos autovectores.}
}
\]

\section{Ejercicio 9}
\subsection{Solución}

Para ambas matrices:

\[
P(\lambda)
=
(4-\lambda)(-2-\lambda)^2.
\]

Para $A$:

\[
S_A(4)=\Span\{(1,-1,-2)\},
\]

\[
S_A(-2)=\Span\{(1,0,-1)\}.
\]

Como:

\[
\dim S_A(-2)=1<2,
\]

\[
\boxed{
A\text{ no es diagonalizable.}
}
\]

Para $B$:

\[
S_B(4)=\Span\{(1,1,2)\},
\]

\[
S_B(-2)
=
\Span\{(1,1,0),(0,1,1)\}.
\]

Por tanto:

\[
\boxed{
B\text{ es diagonalizable.}
}
\]

Una matriz de paso es:

\[
\boxed{
P=
\begin{pmatrix}
1&0&1\\
1&1&1\\
0&1&2
\end{pmatrix}
}
\]

y:

\[
\boxed{
D=
\begin{pmatrix}
-2&0&0\\
0&-2&0\\
0&0&4
\end{pmatrix}.
}
\]

\section{Ejercicio 10}
\subsection{Solución}

El polinomio característico es:

\[
P(\lambda)
=
(2-\lambda)(-1-\lambda)\lambda^2.
\]

Para $\lambda=0$:

\[
S(0)
=
\{(z,-z,z,0):z\in\R\}.
\]

Luego:

\[
\dim S(0)=1<2.
\]

Por tanto:

\[
\boxed{
A\text{ no es diagonalizable.}
}
\]

\section{Ejercicio 11}
\subsection{Solución}

\[
P(\lambda)
=
(3-\lambda)(\lambda^2-2\lambda+3).
\]

Los autovalores son:

\[
3,\qquad
1+i\sqrt2,\qquad
1-i\sqrt2.
\]

Por tanto:

\[
\boxed{
A\text{ no es diagonalizable en }\R.
}
\]

Como en $\C$ tiene tres autovalores distintos:

\[
\boxed{
A\text{ sí es diagonalizable en }\C.
}
\]

\section{Ejercicio 12}
\subsection{Solución}

\[
P(\lambda)=\lambda^4.
\]

El único autovalor es:

\[
\lambda=0
\]

con multiplicidad algebraica $4$.

Además:

\[
S(0)
=
\Ker A
=
\Span\{(1,0,0,0)\},
\]

de modo que:

\[
\dim S(0)=1<4.
\]

Por tanto:

\[
\boxed{
A\text{ no es diagonalizable.}
}
\]

\section{Ejercicio 13}
\subsection{Solución}

La matriz es:

\[
A=
\begin{pmatrix}
3&0&0\\
0&-1&a\\
3&0&b
\end{pmatrix}.
\]

El polinomio característico es:

\[
(3-\lambda)(-1-\lambda)(b-\lambda).
\]

Si:

\[
b\neq-1,3,
\]

hay tres autovalores distintos y:

\[
A\text{ es diagonalizable.}
\]

Si:

\[
b=-1,
\]

entonces $\lambda=-1$ es doble. Se tiene:

\[
\dim S(-1)=2
\iff
a=0.
\]

Si:

\[
b=3,
\]

el autovalor $3$ es doble, pero:

\[
\dim S(3)=1.
\]

Por tanto:

\[
\boxed{
A\text{ es diagonalizable}
\iff
b\neq3
\text{ y }
(b\neq-1\text{ o }a=0).
}
\]

\section{Ejercicio 14}
\subsection{Solución}

Las dimensiones son:

\[
\dim\Ker(f-I)=1,
\]

\[
\dim\Ker(f-I)^2=2,
\]

\[
\dim\Ker(f-I)^3=3.
\]

Por tanto, para $\lambda=1$ existe un bloque de Jordan de orden $3$.

Además:

\[
\Ker f=\Span\{(0,0,0,1)\},
\]

por lo que $\lambda=0$ aporta un bloque de orden $1$.

Una base de Jordan es:

\[
v_1=(1,0,0,1),
\]

\[
v_2=(1,1,0,0),
\]

\[
v_3=(1,0,-1,0),
\]

\[
v_4=(0,0,0,1).
\]

Así:

\[
\boxed{
J=
\begin{pmatrix}
1&0&0&0\\
1&1&0&0\\
0&1&1&0\\
0&0&0&0
\end{pmatrix}.
}
\]

La matriz de paso es:

\[
P=
\begin{pmatrix}
1&1&1&0\\
0&1&0&0\\
0&0&-1&0\\
1&0&0&1
\end{pmatrix}.
\]

Entonces:

\[
A=PJP^{-1}.
\]

Se obtiene:

\[
\boxed{
A=
\begin{pmatrix}
2&0&1&0\\
1&0&1&0\\
0&-1&1&0\\
1&-1&1&0
\end{pmatrix}.
}
\]

Por tanto:

\[
\boxed{
f(x_1,x_2,x_3,x_4)
=
(2x_1+x_3,\;
x_1+x_3,\;
-x_2+x_3,\;
x_1-x_2+x_3).
}
\]

Los autovalores de $A$ son:

\[
1,1,1,0.
\]

Los de:

\[
A-2I
\]

son:

\[
-1,-1,-1,-2.
\]

Por tanto, los de $(A-2I)^{-2}$ son:

\[
\boxed{
1\text{ con multiplicidad }3,
\qquad
\frac14\text{ con multiplicidad }1.
}
\]

\section{Ejercicio 15}
\subsection{Solución}

El polinomio característico es:

\[
P(\lambda)
=
\lambda^2-2(t+2)\lambda+4t+13.
\]

Sus raíces son:

\[
\lambda=t+2\pm\sqrt{t^2-9}.
\]

Si:

\[
|t|<3,
\]

los autovalores no son reales.

Si:

\[
|t|>3,
\]

los autovalores son reales y distintos.

Para:

\[
t=\pm3,
\]

el autovalor es doble y el subespacio propio tiene dimensión $1$.

Por tanto:

\[
\boxed{
A\text{ es diagonalizable en }\R
\iff
|t|>3.
}
\]

\section{Ejercicio 16}
\subsection{Solución}

El polinomio característico es:

\[
P(\lambda)
=
(2-\lambda)^3(-2-\lambda).
\]

Los subespacios propios son:

\[
S(2)
=
\Span\{
(1,1,0,0),
(1,0,1,0),
(1,0,0,1)
\},
\]

\[
S(-2)
=
\Span\{(-1,1,1,1)\}.
\]

Por tanto:

\[
\boxed{
A\text{ es diagonalizable.}
}
\]

Una matriz de paso es:

\[
\boxed{
P=
\begin{pmatrix}
1&1&1&-1\\
1&0&0&1\\
0&1&0&1\\
0&0&1&1
\end{pmatrix}
}
\]

con:

\[
\boxed{
D=
\diag(2,2,2,-2).
}
\]

Para $A^2$:

\[
\boxed{
A^2
=
P\diag(4,4,4,4)P^{-1}.
}
\]

Para $A^{-1}$:

\[
\boxed{
A^{-1}
=
P
\diag
\left(
\frac12,\frac12,\frac12,-\frac12
\right)
P^{-1}.
}
\]

\section{Ejercicio 17}
\subsection{Solución}

El polinomio característico es:

\[
P(\lambda)
=
(1-\lambda)(\lambda^2-4\lambda+5).
\]

Los autovalores son:

\[
1,\qquad2+i,\qquad2-i.
\]

Por tanto:

\[
\boxed{
M\text{ no es diagonalizable en }\R.
}
\]

En $\C$ son tres autovalores distintos.

Bases de los subespacios propios:

\[
S(1)=\Span\{(1,2,1)\},
\]

\[
S(2+i)=\Span\{(i,1+i,2)\},
\]

\[
S(2-i)=\Span\{(i,i-1,-2)\}.
\]

Así:

\[
\boxed{
P=
\begin{pmatrix}
1&i&i\\
2&1+i&i-1\\
1&2&-2
\end{pmatrix}
}
\]

y:

\[
\boxed{
D=
\diag(1,2+i,2-i).
}
\]

\section{Ejercicio 18}
\subsection{Solución}

Como:

\[
\rg A=1,
\]

\[
\dim\Ker A=5.
\]

Por tanto:

\[
\dim S(0)=5.
\]

Así, la multiplicidad algebraica de $\lambda=0$ sólo puede ser $5$ o $6$.

Si es $6$, el único autovalor es $0$ y:

\[
g_0=5<6,
\]

por lo que no es diagonalizable. Su forma de Jordan tiene un bloque $J_2(0)$ y cuatro bloques $J_1(0)$.

Si la multiplicidad algebraica de $0$ es $5$, existe otro autovalor $k\neq0$, simple. Entonces hay seis bloques $1\times1$:

\[
\boxed{
J=
\diag(0,0,0,0,0,k),
}
\]

y la matriz es diagonalizable.

\section{Ejercicio 19}
\subsection{Solución}

Usamos:

\[
\dim E_r(\lambda)
=
6-\rg(A-\lambda I)^r.
\]

En el caso a):

\[
\dim E_1=2,\qquad
\dim E_2=4,\qquad
\dim E_3=5,\qquad
\dim E_4=6.
\]

Las diferencias:

\[
2,\quad2,\quad1,\quad1
\]

corresponden a dos cadenas de Jordan, de longitudes $4$ y $2$.

Por tanto:

\[
\boxed{
J=
J_4(\lambda)\oplus J_2(\lambda).
}
\]

En el caso b):

\[
\dim E_1=3,\qquad
\dim E_2=4,\qquad
\dim E_3=5,\qquad
\dim E_4=6.
\]

Las diferencias son:

\[
3,\quad1,\quad1,\quad1.
\]

Por tanto hay tres bloques, de tamaños $4,1,1$:

\[
\boxed{
J=
J_4(\lambda)
\oplus
J_1(\lambda)
\oplus
J_1(\lambda).
}
\]

Ambos casos son posibles.

\section{Ejercicio 20}
\subsection{Solución}

Sean:

\[
v_1=(1,0,1),
\quad
v_2=(-1,1,0),
\quad
v_3=(0,1,-1).
\]

Imponiendo:

\[
Av_i=\lambda_i v_i
\]

se obtiene:

\[
\lambda_1=6,
\qquad
\lambda_2=-4,
\qquad
\lambda_3=0.
\]

La matriz resultante es:

\[
\boxed{
A=
\begin{pmatrix}
1&5&5\\
2&-2&-2\\
3&3&3
\end{pmatrix}.
}
\]

Por tanto:

\[
\boxed{
a=5,\quad b=-2,\quad c=3,\quad d=5,\quad e=-2,\quad f=3.
}
\]

Para admitir diagonalización ortogonal, una matriz real debe ser simétrica. Esto exigiría, entre otras condiciones:

\[
a=2.
\]

Pero aquí:

\[
a=5.
\]

Por tanto:

\[
\boxed{
\text{Con las condiciones anteriores no admite diagonalización ortogonal.}
}
\]

\section{Ejercicio 21}
\subsection{Solución}

La evolución viene dada por:

\[
X_{k+1}
=
AX_k,
\]

con:

\[
A=
\frac13
\begin{pmatrix}
0&1&1\\
2&2&0\\
1&0&2
\end{pmatrix}.
\]

Sea:

\[
M=
\begin{pmatrix}
0&1&1\\
2&2&0\\
1&0&2
\end{pmatrix}.
\]

Sus autovalores son:

\[
2,\qquad3,\qquad-1.
\]

El autovalor dominante de $A=\frac13M$ es:

\[
1.
\]

Un autovector asociado es:

\[
(1,2,1).
\]

La población total se conserva:

\[
60+200+300=560.
\]

En el estado estacionario:

\[
(A,B,C)=k(1,2,1).
\]

Entonces:

\[
4k=560
\]

y:

\[
k=140.
\]

Por tanto:

\[
\boxed{
(A,B,C)\longrightarrow(140,280,140).
}
\]

\section{Ejercicio 22}
\subsection{Solución}

El polinomio característico es:

\[
(5-\lambda)(-1-\lambda)(a-\lambda).
\]

Si:

\[
a\neq-1,5,
\]

hay tres autovalores distintos y $A$ es diagonalizable.

Si:

\[
a=-1,
\]

el autovalor $-1$ es doble y:

\[
\dim S(-1)=2
\iff
b=0.
\]

Si:

\[
a=5,
\]

el autovalor $5$ es doble y su multiplicidad geométrica es $1$.

Por tanto:

\[
\boxed{
A\text{ es diagonalizable}
\iff
a\neq5
\text{ y }
(a\neq-1\text{ o }b=0).
}
\]

Para:

\[
a=-1,\qquad b=-1,
\]

los autovalores son:

\[
-1\text{ doble},\qquad5\text{ simple}.
\]

La forma de Jordan es:

\[
\boxed{
J=
\begin{pmatrix}
-1&0&0\\
1&-1&0\\
0&0&5
\end{pmatrix}.
}
\]

Para Cayley--Hamilton:

\[
p(\lambda)
=
-(\lambda+1)^2(\lambda-5)
=
-\lambda^3+3\lambda^2+9\lambda+5.
\]

El resto de dividir $\lambda^9$ por $p(\lambda)$ es:

\[
r(\lambda)
=
54252\lambda^2
+
108513\lambda
+
54260.
\]

Por tanto:

\[
\boxed{
A^9
=
54252A^2
+
108513A
+
54260I.
}
\]

\section{Ejercicio 23}
\subsection{Solución}

La matriz tiene rango:

\[
\rg M=1.
\]

Por tanto:

\[
\dim S(0)
=
3-\rg M
=
2.
\]

Además:

\[
\tr M=3,
\]

por lo que los autovalores son:

\[
0,\quad0,\quad3.
\]

Para $\lambda=0$:

\[
S(0)
=
\Span\{(-2,1,0),(-3,0,1)\}.
\]

Para $\lambda=3$:

\[
S(3)
=
\Span\{(0,3,2)\}.
\]

Tenemos tres autovectores linealmente independientes, luego:

\[
\boxed{
M\text{ es diagonalizable.}
}
\]

Una base de autovectores de $\mathcal P_2(x)$ es:

\[
\boxed{
\{-2+x,\;-3+x^2,\;3x+2x^2\}.
}
\]

\section{Ejercicio 24}
\subsection{Solución}

El polinomio característico es:

\[
p(\lambda)
=
-(\lambda+1)^2(\lambda-2).
\]

Por tanto:

\[
\lambda=-1
\]

tiene multiplicidad algebraica $2$ y:

\[
\lambda=2
\]

es simple.

Se obtiene:

\[
S(-1)
=
\Span\{(-3,1,0),(-9,0,1)\},
\]

\[
S(2)
=
\Span\{(9,3,1)\}.
\]

Así:

\[
\boxed{
M\text{ es diagonalizable.}
}
\]

Una base de autovectores de $\mathcal P_2(x)$ es:

\[
\boxed{
\{-3+x,\;-9+x^2,\;9+3x+x^2\}.
}
\]

Por Cayley--Hamilton:

\[
p(A)
=
-A^3+3A+2I=0.
\]

Entonces:

\[
2I=A(A^2-3I),
\]

y:

\[
\boxed{
M^{-1}
=
\frac12(M^2-3I).
}
\]

Si:

\[
P=
\begin{pmatrix}
-3&-9&9\\
1&0&3\\
0&1&1
\end{pmatrix},
\qquad
D=
\diag(-1,-1,2),
\]

entonces:

\[
\boxed{
M^p
=
PD^pP^{-1}.
}
\]

\section{Ejercicio 25}
\subsection{Solución}

El polinomio característico es:

\[
p(\lambda)
=
(\lambda+1)^2(\lambda-2)^2
\]

o:

\[
p(\lambda)
=
\lambda^4-2\lambda^3-3\lambda^2+4\lambda+4.
\]

Por Cayley--Hamilton:

\[
A^4-2A^3-3A^2+4A+4I=0.
\]

Como $A$ es regular:

\[
I
=
\frac14
A(-A^3+2A^2+3A-4I).
\]

Por tanto:

\[
\boxed{
A^{-1}
=
\frac14
\left(
-A^3+2A^2+3A-4I
\right).
}
\]

\section{Ejercicio 26}
\subsection{Solución}

El polinomio característico es:

\[
p(\lambda)
=
(\lambda-1)(\lambda-6).
\]

Los autovectores pueden tomarse:

\[
v_1=(2,-3)
\]

para $\lambda=1$ y:

\[
v_2=(1,1)
\]

para $\lambda=6$.

Por tanto:

\[
P=
\begin{pmatrix}
2&1\\
-3&1
\end{pmatrix},
\qquad
D=
\begin{pmatrix}
1&0\\
0&6
\end{pmatrix}.
\]

Así:

\[
\boxed{
M^n
=
PD^nP^{-1}.
}
\]

Por Cayley--Hamilton:

\[
M^2-7M+6I=0.
\]

Buscamos:

\[
M^n=aM+bI.
\]

Imponiendo los autovalores:

\[
1=a+b,
\]

\[
6^n=6a+b.
\]

Por tanto:

\[
a=\frac{6^n-1}{5},
\qquad
b=\frac{6-6^n}{5}.
\]

Así:

\[
\boxed{
M^n
=
\frac{6^n-1}{5}M
+
\frac{6-6^n}{5}I.
}
\]

\section{Ejercicio 27}
\subsection{Solución}

\[
\tr A=2,
\qquad
\det A=1.
\]

Por tanto:

\[
p(\lambda)
=
(\lambda-1)^2.
\]

Cayley--Hamilton da:

\[
(A-I)^2=0.
\]

Se obtiene:

\[
A^n
=
nA+(1-n)I.
\]

Por tanto:

\[
\frac1nA^n
=
A+
\left(
\frac1n-1
\right)I.
\]

Así:

\[
\boxed{
\lim_{n\to\infty}\frac1nA^n
=
A-I
=
\begin{pmatrix}
-15&25\\
-9&15
\end{pmatrix}.
}
\]

\section{Ejercicio 28}
\subsection{Solución}

El único autovalor es:

\[
\lambda=3
\]

con multiplicidad algebraica $5$.

Las dimensiones son:

\[
\dim E_1(3)=2,
\qquad
\dim E_2(3)=4,
\qquad
\dim E_3(3)=5.
\]

Por tanto hay dos bloques de Jordan, de tamaños $3$ y $2$.

Una base de Jordan es:

\[
v_1=(1,0,0,0,0),
\]

\[
v_2=(1,-1,1,1,0),
\]

\[
v_3=(0,0,1,0,0),
\]

\[
v_4=(0,0,0,1,1),
\]

\[
v_5=(0,1,0,1,0).
\]

Así:

\[
\boxed{
J=
\begin{pmatrix}
3&0&0&0&0\\
1&3&0&0&0\\
0&1&3&0&0\\
0&0&0&3&0\\
0&0&0&1&3
\end{pmatrix}.
}
\]

Una matriz de paso es:

\[
\boxed{
P=
\begin{pmatrix}
1&1&0&0&0\\
0&-1&0&0&1\\
0&1&1&0&0\\
0&1&0&1&1\\
0&0&0&1&0
\end{pmatrix}.
}
\]

\section{Ejercicio 29}
\subsection{Solución}

Los autovalores son:

\[
2
\quad\text{y}\quad
-2,
\]

ambos con multiplicidad algebraica $2$.

Para ambos:

\[
g_\lambda=1.
\]

Por tanto existe un bloque $2\times2$ para cada autovalor.

Una base de Jordan es:

\[
v_1=(0,1,0,1),
\qquad
v_2=(0,0,1,0),
\]

\[
v_3=(1,1,0,0),
\qquad
v_4=(1,0,1,0).
\]

Así:

\[
\boxed{
J=
\begin{pmatrix}
2&0&0&0\\
1&2&0&0\\
0&0&-2&0\\
0&0&1&-2
\end{pmatrix}.
}
\]

y:

\[
\boxed{
P=
\begin{pmatrix}
0&0&1&1\\
1&0&1&0\\
0&1&0&1\\
1&0&0&0
\end{pmatrix}.
}
\]

\section{Ejercicio 30}
\subsection{Solución}

El polinomio característico es:

\[
p(\lambda)
=
(1-\lambda)^2(1+\alpha-\lambda).
\]

Si:

\[
\alpha\neq0,
\]

los autovalores son:

\[
1\text{ doble},
\qquad
1+\alpha\text{ simple}.
\]

Para $\lambda=1$:

\[
\dim S(1)=1,
\]

por lo que aparece un bloque $2\times2$.

Una base de Jordan es:

\[
v_1=(1,0,0),
\]

\[
v_2=(0,1,1),
\]

\[
v_3=(\alpha,1,1+\alpha^2).
\]

Por tanto:

\[
\boxed{
J=
\begin{pmatrix}
1&0&0\\
1&1&0\\
0&0&1+\alpha
\end{pmatrix},
\qquad
\alpha\neq0.
}
\]

Si:

\[
\alpha=0,
\]

el único autovalor es $1$ con multiplicidad $3$ y:

\[
\dim E_1=1,\qquad
\dim E_2=2,\qquad
\dim E_3=3.
\]

Por tanto:

\[
\boxed{
J=
\begin{pmatrix}
1&0&0\\
1&1&0\\
0&1&1
\end{pmatrix},
\qquad
\alpha=0.
}
\]

\section{Ejercicio 31}
\subsection{Solución}

El único autovalor es:

\[
\lambda=1
\]

con multiplicidad algebraica $4$.

Si:

\[
k\neq0,
\]

las dimensiones de los núcleos sucesivos conducen a bloques de tamaños $3$ y $1$.

Por tanto:

\[
\boxed{
J=
\begin{pmatrix}
1&0&0&0\\
1&1&0&0\\
0&1&1&0\\
0&0&0&1
\end{pmatrix}.
}
\]

Una matriz de paso proporcionada por las cadenas de Jordan es:

\[
\boxed{
P=
\begin{pmatrix}
0&0&k&0\\
0&k&0&0\\
0&k&k&1\\
1&0&-k&1
\end{pmatrix}.
}
\]

Si:

\[
k=0,
\]

los bloques tienen tamaños:

\[
2,1,1.
\]

Por tanto:

\[
\boxed{
J=
\begin{pmatrix}
1&0&0&0\\
1&1&0&0\\
0&0&1&0\\
0&0&0&1
\end{pmatrix}.
}
\]

Una matriz de paso es:

\[
\boxed{
P=
\begin{pmatrix}
0&1&1&0\\
1&0&0&0\\
0&1&0&1\\
1&-1&0&0
\end{pmatrix}.
}
\]

\section{Ejercicio 32}
\subsection{Solución}

Para:

\[
\lambda=1
\]

la multiplicidad algebraica es $2$ y la geométrica $1$, por lo que aparece:

\[
J_2(1).
\]

Para:

\[
2+i
\]

y su conjugado:

\[
2-i,
\]

la multiplicidad algebraica es $2$ y la geométrica $1$.

El bloque real básico es:

\[
C=
\begin{pmatrix}
2&1\\
-1&2
\end{pmatrix}.
\]

Como la multiplicidad es $2$, aparece un bloque real de Jordan de orden $2$.

Por tanto, una forma real de Jordan es:

\[
\boxed{
J_R=
\begin{pmatrix}
1&0&0&0&0&0\\
1&1&0&0&0&0\\
0&0&2&1&0&0\\
0&0&-1&2&0&0\\
0&0&1&0&2&1\\
0&0&0&1&-1&2
\end{pmatrix}.
}
\]

\section{Ejercicio 33}
\subsection{Solución}

El polinomio característico es:

\[
p(\lambda)
=
(\lambda^2+1)^2.
\]

Los autovalores complejos son:

\[
i,\qquad-i,
\]

ambos de multiplicidad algebraica $2$.

Para $\lambda=i$ puede tomarse la cadena:

\[
v_1=(0,0,1,-1-i),
\]

\[
v_2=(M-iI)v_1=(-i,1+i,0,0).
\]

Escribimos:

\[
v_1=u_1+iw_1,
\qquad
v_2=u_2+iw_2,
\]

con:

\[
u_1=(0,0,1,-1),
\qquad
w_1=(0,0,0,-1),
\]

\[
u_2=(0,1,0,0),
\qquad
w_2=(-1,1,0,0).
\]

Una base real es:

\[
\boxed{
B'=
\{u_1,w_1,u_2,w_2\}.
}
\]

En ella:

\[
\boxed{
J_R=
\begin{pmatrix}
0&1&0&0\\
-1&0&0&0\\
1&0&0&1\\
0&1&-1&0
\end{pmatrix}.
}
\]

\section{Ejercicio 34}
\subsection{Solución}

Para:

\[
A=
\begin{pmatrix}
-2&0&0\\
0&2&0\\
0&1&2
\end{pmatrix},
\]

el bloque asociado a $2$ tiene parte nilpotente de orden $2$.

Por tanto:

\[
\boxed{
e^A
=
\begin{pmatrix}
e^{-2}&0&0\\
0&e^2&0\\
0&e^2&e^2
\end{pmatrix}.
}
\]

Más generalmente:

\[
\boxed{
e^{At}
=
\begin{pmatrix}
e^{-2t}&0&0\\
0&e^{2t}&0\\
0&te^{2t}&e^{2t}
\end{pmatrix}.
}
\]

Para:

\[
B=
\begin{pmatrix}
3&0&0\\
1&3&0\\
0&1&3
\end{pmatrix},
\]

se tiene un único bloque de Jordan de orden $3$.

Así:

\[
\boxed{
e^B
=
e^3
\begin{pmatrix}
1&0&0\\
1&1&0\\
\frac12&1&1
\end{pmatrix}.
}
\]

y:

\[
\boxed{
e^{Bt}
=
e^{3t}
\begin{pmatrix}
1&0&0\\
t&1&0\\
\frac{t^2}{2}&t&1
\end{pmatrix}.
}
\]

\section{Ejercicio 35}
\subsection{Solución}

Como:

\[
M=\lambda I,
\]

\[
M^k=\lambda^kI.
\]

Entonces:

\[
e^M
=
\sum_{k=0}^\infty
\frac{\lambda^k}{k!}I.
\]

Por tanto:

\[
\boxed{
e^M=e^\lambda I_n.
}
\]

\section{Ejercicio 36}
\subsection{Solución}

La matriz es nilpotente y:

\[
A^n=0.
\]

Por tanto la serie termina:

\[
\boxed{
e^A
=
I+A+\frac{A^2}{2!}
+\cdots+
\frac{A^{n-1}}{(n-1)!}.
}
\]

En forma matricial:

\[
\boxed{
e^A=
\begin{pmatrix}
1&0&0&\cdots&0\\
1&1&0&\cdots&0\\
\frac1{2!}&1&1&\ddots&\vdots\\
\vdots&\ddots&\ddots&\ddots&0\\
\frac1{(n-1)!}&\cdots&\frac1{2!}&1&1
\end{pmatrix}.
}
\]

\section{Ejercicio 37}
\subsection{Solución}

Escribimos:

\[
A=
aI+
\begin{pmatrix}
0&b\\
-b&0
\end{pmatrix}.
\]

Las dos matrices conmutan, luego:

\[
e^A
=
e^{aI}
e^{
\left(
\begin{smallmatrix}
0&b\\
-b&0
\end{smallmatrix}
\right)
}.
\]

La segunda exponencial es una rotación:

\[
\boxed{
e^A
=
e^a
\begin{pmatrix}
\cos b&\sen b\\
-\sen b&\cos b
\end{pmatrix}.
}
\]

\section{Ejercicio 38}
\subsection{Solución}

El polinomio característico es:

\[
p(\lambda)
=
(1-\lambda)^2
\left(
\frac12-\lambda
\right).
\]

Desarrollando:

\[
p(\lambda)
=
-\lambda^3
+
\frac52\lambda^2
-
2\lambda
+
\frac12.
\]

Por Cayley--Hamilton:

\[
-A^3
+
\frac52A^2
-
2A
+
\frac12I
=
0.
\]

Multiplicando por $2$:

\[
-2A^3+5A^2-4A+I=0.
\]

Multiplicando por $A$:

\[
-2A^4+5A^3-4A^2+A=0.
\]

Restando adecuadamente ambas identidades:

\[
\boxed{
2A^4-7A^3+9A^2-5A+I=0.
}
\]

\section{Ejercicio 39}
\subsection{Solución}

Como $B$ tiene tres autovalores distintos:

\[
0,1,2,
\]

es diagonalizable.

Su rango coincide con el de:

\[
D=\diag(0,1,2).
\]

Por tanto:

\[
\boxed{
\rg B=2.
}
\]

Además:

\[
\det B=0.
\]

Entonces:

\[
\det(B^tB)
=
\det(B^t)\det(B)
=
(\det B)^2.
\]

Por tanto:

\[
\boxed{
\det(B^tB)=0.
}
\]

Los autovalores de $B^tB$ no pueden deducirse únicamente a partir de los autovalores de $B$:

\[
\boxed{
\text{No hay información suficiente para el apartado c).}
}
\]

Para $B^2$ los autovalores son:

\[
0,1,4.
\]

Para:

\[
B^2+I
\]

son:

\[
1,2,5.
\]

Por tanto, para la inversa:

\[
\boxed{
(B^2+I)^{-1}:
\quad
1,\frac12,\frac15.
}
\]

\end{document}
